\section{Markov Benchmark Details}\label{app:min-markov}

\subsection{Data-Generating Process}

Each synthetic document is a fixed-length token sequence generated by a
hidden Markov chain over $K$ registers. The chain picks a register; the
register emits an observed token from its own palette of synonyms---
register $A$ emits \emph{grief}, \emph{mourn}, \emph{tears}; register $B$
emits \emph{plot}, \emph{spy}, \emph{ploy}; and so on. The learner sees
tokens, not hidden registers. The oracle is the number of register
transitions in the hidden sequence.

\begin{figure}[t]
\centering
\includegraphics[width=\linewidth]{markov_scene_arc_hamlet.pdf}
\caption{\textbf{Boundary correction on Hamlet's Act 3 through
Ophelia's funeral.} Every scene from 3.1 to 5.1 (twelve in total) is
colored by its dominant register from the five-register set
Mourning, Intrigue, Madness, Action, Political. Three same-register
runs are visible in the strip: 3.2--3.3 (Intrigue), 4.1--4.2--4.3
(Political), and 4.4--4.5 (Mourning). Leaves $L_1,\ldots,L_4$ group
three consecutive scenes and each leaf summary $S_i = g(L_i)$ stores
$(\text{first register}, \text{last register}, \text{count})$; each
internal node adds $c_L + c_R + \One\{L.\mathrm{last}\neq
R.\mathrm{first}\}$ and the root returns the exact count of register
shifts across the arc ($=7$). Two canonical scene-boundary flips
sit between leaves: Claudius-at-prayer to Polonius-killed
(Intrigue$\to$Action, 3.3$\to$3.4) and the foil plot to Ophelia's
funeral (Intrigue$\to$Mourning, 4.7$\to$5.1). At the root merge the
boundary correction correctly does \emph{not} fire
($E\text{\,vs\,}E$), because the Political run that ends $L_2$
continues into $L_3$; the tree sees this because the leaf summaries
retain their first and last register. Source:
\texttt{paper/figures/markov\_scene\_arc\_hamlet.tex}.}
\label{fig:min-markov-registers}
\end{figure}

\subsection{Sufficient State}

For any span, the sufficient state is
\[
(\mathrm{count},\mathrm{first},\mathrm{last}).
\]
The count records internal transitions. The first and last fields carry the
boundary information needed by parent merges. The merge rule is
\[
(c_L,a_L,b_L)\otimes(c_R,a_R,b_R)
=
(c_L+c_R+\One\{b_L\ne a_R\},a_L,b_R).
\]
This state is associative and ordered. A leaf-only or unordered aggregate loses
boundary transitions and cannot be repaired at the root.

\subsection{Why This Synthetic Setting}

The synthetic chain exposes the local-law issue in the smallest setting
that contains it: a parent merge needs the right boundary of the left
child and the left boundary of the right child. When either child drops
that boundary state, the root count shifts even when each child summary
looks locally adequate---the boundary-correction pattern shown in
Figure~\ref{fig:min-markov-registers}. The Hamlet labels in the figure
are a reader-facing skin on the synthetic registers; every number in the
empirical results below is computed on the synthetic palette.

\subsection{Empirical Role}

The Markov experiment is the mechanism check for local supervision. A flat FNO
learns from root labels. The tree model learns from root labels plus local
leaf/merge labels. The retained result is that tree supervision matches or
improves on the flat baseline across the budget range and remains competitive
when only $10\%$ of documents receive root labels.
Figure~\ref{fig:min-markov-leaf-size} reports these per-cell gains across
leaf size and the fraction of documents receiving root labels.
This supports the main
granularity claim in a setting where the sufficient state is known and
node-level labels are deliberately observed. The Manifesto/Benoit experiment in
the main text is different: it observes full-document root targets over the
sampled corpus and reserves node-level observations for the stronger
quasi-sentence audit path.

\begin{figure}[t]
\centering
\begin{subfigure}[b]{0.48\textwidth}
    \centering
    \includegraphics[width=\linewidth]{markov_leaf_size_fixed_recoverable.pdf}
    \caption{Recoverable scope.}
    \label{fig:min-markov-leaf-size-recoverable}
\end{subfigure}
\hfill
\begin{subfigure}[b]{0.48\textwidth}
    \centering
    \includegraphics[width=\linewidth]{markov_leaf_size_fixed_structural.pdf}
    \caption{Structural scope.}
    \label{fig:min-markov-leaf-size-structural}
\end{subfigure}
\caption{\textbf{Tree-supervision gain over the flat FNO baseline across
leaf size and supervision budget.} Recoverable (left) and structural
(right) scopes on the synthetic Markov chain. The bolded colorbar tick
marks the average flat-FNO root MAE for the panel, so cells above the
tick improve on the flat baseline.}
\label{fig:min-markov-leaf-size}
\end{figure}
